\subsection{An explicit continuous class-three lifting obstruction}

\begin{statement}{A.196}{Proved}
{Every finite class-three quotient lifts, while the continuous class-three quotient does not}
Use the field, elements, and fixed compatible Kummer cocycles of
A.195:
\[
 \begin{gathered}
 k=\Q(\alpha),\quad\alpha^3=3,\quad d=649,\quad
 A=d^8\alpha,\quad U=\alpha^2-2,\quad B=U^3,\\
 a=\kappa_3(A),\qquad b=\kappa_3(B).
 \end{gathered}
\tag{A.196.1}
\]
Here \(\kappa_3\) denotes the integral three-primary Kummer
cocycle, and its reductions are always taken from this same
compatible root system.  There is an actual arithmetic quotient
of the tripod fundamental group over \(k\) with geometric group
\[
 \mathcal Q=(R y)\rtimes\Z_3x,
 \qquad R=\Z_3[T]/T^3,\quad X=1+T,
 \qquad xvx^{-1}=Xv.
\tag{A.196.2}
\]
Its lower central series satisfies
\(\Gamma_2\mathcal Q=TR\),
\(\Gamma_3\mathcal Q=T^2R\ne0\), and
\(\Gamma_4\mathcal Q=0\).  Its derived subgroup is abelian.

The abelian arithmetic section specified by \((a,b)\) has a
continuous global lift to \(\mathcal Q/\Gamma_3\mathcal Q\),
and has continuous local lifts to \(\mathcal Q\) at every place
of \(k\).  It has no continuous global lift to \(\mathcal Q\).
The obstruction to the last lift is exactly the A.195 class
\(\mathfrak o_2\), and this class has order three.

For every \(c\ge1\), it nevertheless has a continuous global
lift to the arithmetic quotient with finite geometric group
\[
 \mathcal Q_c=
 \bigl((\Z/3^c)[T]/T^3\bigr)\rtimes\Z/3^c,
 \qquad |\mathcal Q_c|=3^{4c},
 \qquad \operatorname{class}(\mathcal Q_c)=3.
\tag{A.196.3}
\]
These form a cofinal system with \(\mathcal Q=\varprojlim_c
\mathcal Q_c\).  There is no compatible system of these global
lifts, even after allowing geometric conjugacy at each level.

Put \(S=\{v:v\mid2\cdot3\cdot11\cdot59\}\cup S_\infty\).
For every finite \(S'\supset S\), there is an integer
\(c(S')\) such that no lift at any level \(c\ge c(S')\)
has its geometric cocycle factoring through \(G_{k,S'}\).
For \(S'=S\), the first failing level is exactly \(c=2\).
All assertions refer to the actual arithmetic pullback maps
specified below.  The local data in this statement are exactly
the continuous norm-derived quotient sections of A.181.  At the
three-adic place, their prescribed abelian frame modulo nine
does not occur for an ordinary point or a cuspidal section;
the literal ninth-power cover and all three cusp images are
checked in the proof.
\end{statement}

\begin{proof}
\emph{The quotient of the original arithmetic fundamental group.}
Let \(\Pi\) be the geometric free pro-three group of the
tripod, with the tangential generators \(x,y\) used in A.181.
Set \(N=\ker(\Pi\to\Z_3(1))\), where \(x\mapsto1\) and
\(y\mapsto0\).  Close all commutator subgroups.  The actual
Schreier and Fox maps of A.77 and A.181 identify
\[
 \Pi/[N,N]=(\Z_3[[T]]y)\rtimes\Z_3x,
 \qquad
 h(X,Y)[x,y]\longmapsto h(X,1)Ty.
\tag{A.196.4}
\]
Thus the kernel from \(\Pi/\Pi''\) is precisely
\((Y-1)D\), in the displayed derived coordinates.  Quotient
the normal module in A.196.4 by \(T^3\).  This gives A.196.2
with its actual quotient homomorphism from \(\Pi\).

The tangential arithmetic action has
\(g(x)=x^{\chi(g)}\) and
\(g(y)=f_g^{-1}y^{\chi(g)}f_g\), with
\(f_g\in\Pi'\subset N\).  In \(N/[N,N]\) the images of
\(f_g\) and \(y\) commute.  Consequently this original action
descends to
\[
 g(v,a_0)=
 \bigl(\chi(g)v(X^{\chi(g)}-1),\,\chi(g)a_0\bigr)
 \quad(v\in R,\ a_0\in\Z_3).
\tag{A.196.5}
\]
The ideal \((T^3)\) is preserved because
\(X^{\chi(g)}-1\) is \(T\) times a unit.  No change of the
arithmetic action is made in forming this quotient.

The semidirect product law is
\((v,a_0)(v',a'_0)=(v+X^{a_0}v',a_0+a'_0)\).
Here \([x,y]=Ty\); conjugating a normal-module element by
\(x\) and subtracting it multiplies it by \(T\).  The normal
module is abelian, so these computations give
\[
 \Gamma_j\mathcal Q=T^{j-1}R\quad(j\ge2).
\tag{A.196.6}
\]
The same calculation after reduction modulo \(3^c\) proves
the exact lower-central class in A.196.3.  Separately,
\(\mathcal Q'=TR\) is abelian, so \(\mathcal Q''=1\).
In particular the maps from \(\Pi/\Gamma_4\Pi\) and
\(\Pi/\Pi''\) to \(\mathcal Q\) are actual quotient maps.

For an explicit coordinate comparison, write
\(v=b_0+c_0T+e_0T^2\) and set \(z_0=e_0+c_0/2\).
Division by two is integral over \(\Z_3\).  Then
\(\mathcal Q\) is exactly the following matrix group:
\[
 (a_0,b_0,c_0,z_0)\longmapsto
 \begin{pmatrix}
 1&a_0&a_0^2/2&z_0\\
 0&1&a_0&c_0\\
 0&0&1&b_0\\
 0&0&0&1
 \end{pmatrix}.
\tag{A.196.7}
\]
Matrix multiplication gives
\[
\begin{aligned}
 a''_0&=a_0+a'_0,& b''_0&=b_0+b'_0,\\
 c''_0&=c_0+c'_0+a_0b'_0,&
 z''_0&=z_0+z'_0+a_0c'_0+\tfrac12a_0^2b'_0.
\end{aligned}
\tag{A.196.8}
\]
Under A.196.5 the four coordinates have weights
\(\chi,\chi,\chi^2,\chi^3\), respectively.  Equivalently,
the action is conjugation by
\(\operatorname{diag}(\chi^3,\chi^2,\chi,1)\).
After discarding \(z_0\), A.196.8 is the actual Heisenberg
model of \(\Pi/\Gamma_3\Pi\) in A.38.  This proves the
claimed full class-two comparison, including the arithmetic
action.

\emph{The prescribed lifting problem and its class.}
Twist the normal-module action in A.196.5 by the fixed first
coordinate \(a(g)\).  The resulting module is
\(V^{(3)}_A=V_A/T^3V_A\), with
\[
 gv=\chi(g)X^{a(g)}v(X^{\chi(g)}-1),
 \qquad
 0\longrightarrow D_2\xrightarrow{\ T\ }V^{(3)}_A
      \xrightarrow{\epsilon}\Z_3(1)\longrightarrow0.
\tag{A.196.9}
\]
Multiplication by \(T\) is exactly the A.183 coefficient map.
The group law proves that a section with first coordinate
\(a\) is a continuous one-cocycle in \(V^{(3)}_A\), and its
second abelian coordinate is its augmentation.  An augmentation
coboundary can be removed by lifting its zero-cochain through
\(\epsilon\).  Similarly a discrepancy in the first abelian
cocycle can be removed by conjugation using a lift of its
zero-cochain in the \(x\)-coordinate.  Thus both literal
cocycles and their section classes give the same obstruction
\[
 \delta\kappa_3(B)=\mathfrak o_2\in
                 H^2_{\rm cts}(G_k,D_2).
\tag{A.196.10}
\]
A.195 proves that this class is nonzero.  Its quotient in
\(H^2(G_k,\Z_3(2))\) is zero by A.183, so the same exact
sequence modulo \(T^2V_A\) supplies the continuous class-two
lift.  A.181 constructs a compatible local norm family at
each place.  Its actual local Shapiro map gives a continuous
lift to the full module \(V_A\); quotienting by \(T^3\)
gives the asserted local sections of \(\mathcal Q\).

One can also see every individual central class-three defect.
For a class-two lift with third coordinate \(h\), the equation
is
\(\delta h(g,h')=-a(g)\chi(g)b(h')\).
In A.196.7 its next central defect is the actual cocycle
\[
 C_h(g,h')=
 a(g)\chi(g)^2h(h')+
              \tfrac12a(g)^2\chi(g)b(h')
       \quad\text{in }\Z_3(3).
\tag{A.196.11}
\]
The central extension and A.196.8 show that this is a
two-cocycle; alternatively its differential cancels using
the displayed equation for \(h\).  Under
\(\Z_3(3)\xrightarrow{T}D_2\) its class maps to
\(\mathfrak o_2\).  Hence every such global class-two lift
has a nonzero central defect.  A further coordinate \(z\)
would have to satisfy \(\delta z=-C_h\).

\emph{The exact order of the rank-two obstruction.}
We give the coefficient comparison needed to evaluate its
order.  In the normalized basis \(1-T/2,T\) of \(D_2\),
the action is
\[
 g(1-T/2)=\chi(g)^2(1-T/2)+\chi(g)^2a(g)T,
 \qquad gT=\chi(g)^3T.
\tag{A.196.12}
\]
Thus the exact top/bottom sequence is
\(0\to\Z_3(3)\xrightarrow{T}D_2\to\Z_3(2)\to0\).

For clarity, the twist-three version of the A.183 argument
uses the same actual field \(F=k(\zeta_3)\) and
\(X_0=\operatorname{Spec}\mathcal O_F[1/3]\).
A.138 gives zero three-primary class group and precisely one
prime above three.  The Kummer and Brauer sequences therefore
give \(H^2_{\acute et}(X_0,\mu_3)=0\), as proved in A.183.
The sheaf \(\Z/3(3)\) is identified with \(\mu_3\) by the
chosen \(\zeta_3\).  Induction through
\(0\to\Z/3^{c-1}(3)\to\Z/3^c(3)\to\Z/3(3)\to0\)
gives zero second cohomology at every finite coefficient.
Finite first cohomology removes the inverse-limit term, so
the integral second cohomology on \(X_0\) is zero as well.

The actual localization and degree-two group-to-scheme maps
proved in A.183 now inject the integral group cohomology with
three inverted and the extra primes removed into
\[
 \bigoplus_{v\mid2\cdot11\cdot59}
 H^1(\mathbf F_v,\Z_3(2))
 =\bigoplus_{v\mid2\cdot11\cdot59}
               \Z_3/(|\mathbf F_v|^2-1)\Z_3.
\tag{A.196.13}
\]
There are three primes at each rational prime, with residue
sizes \(4,121,3481\), respectively.  Every displayed
\(|\mathbf F_v|-1\) has three-adic valuation one, and
\(|\mathbf F_v|+1\) is prime to three.  Thus every summand
is killed by three.  Restriction followed by corestriction
for \(F/k\) is multiplication by two; this is invertible
on these coefficient groups.  Consequently
\[
             3H^2(G_{k,S},\Z_3(3))=0.
\tag{A.196.14}
\]
The supported class \(\mathfrak o_{S,2}\) has zero bottom
image by A.183.  Exactness of A.196.12 therefore puts it in
the image of the group in A.196.14, so it is killed by three.
It is nonzero by A.195, and the integral inflation map of
A.183 is injective.  This proves that the particular image
\(\mathfrak o_2\) has exact order three.

\emph{The actual finite quotients and their lifts.}
For \(m=3^c\), both \(m\) and \(\binom m2\) are divisible
by \(3^c\).  Hence
\[
             X^{3^c}-1\in(3^c,T^3)\Z_3[[T]].
\tag{A.196.15}
\]
This makes the \(x\)-action in A.196.3 well defined and
proves that the quotient map reduces both the coefficients
and the \(x\)-coordinate modulo \(3^c\).  The arithmetic
action preserves its kernel, by A.196.5.  The resulting
finite group has four independent \(\Z/3^c\) coordinates
in A.196.7, proving its order.  These coordinate kernels
form a neighbourhood basis of the identity in \(\mathcal Q\),
so they are cofinal and their inverse limit is exactly the
original compact group, with the displayed action.

There is also a literal map from the A.181 root-torsor
quotient to each of these groups:
\[
 \bigl(\Z_3[X]/(X^{3^c}-1)\bigr)\rtimes\Z_3
       \longrightarrow\mathcal Q_c,
 \qquad X\longmapsto1+T,
\tag{A.196.16}
\]
with coefficient and \(x\)-coordinate reduction modulo
\(3^c\).  Equation A.196.15 proves that it is defined.
A.181 supplies a completed norm preimage of \(B\) in
\(K_{c+1}\).  Its actual Shapiro cocycle gives a section
of the source arithmetic extension in A.196.16 with
augmentation \(b\).  Its image is the claimed global lift
to \(\mathcal Q_c\).  Equivalently, the finite coefficient
sequence obtained from A.196.9 has obstruction the absolute
image of \(\mathfrak o_2\) in \(H^2(G_k,D_2/3^c)\),
which A.181 proves to be zero.

A compatible sequence of literal finite cocycles would
give a continuous cocycle to \(\mathcal Q=\varprojlim_c
\mathcal Q_c\), contradicting A.196.10.  The same conclusion
holds for compatible geometric-conjugacy classes.  Indeed,
choose representatives inductively and lift a conjugator
from \(\mathcal Q_c\) to \(\mathcal Q_{c+1}\) to make
successive representatives exactly compatible.  Their
abelian cohomology classes are those of \((a,b)\) at every
level.  Here \(H^0(k,\mu_{3^c})=0\), since \(k\) does
not contain \(\mu_3\); the abelian correcting zero-cochains
are consequently unique and compatible.  The limiting
section has the required integral abelian class, which is
impossible by A.196.10.

\emph{Every fixed ramification set eventually fails.}
The action on each \(\mathcal Q_c\), and the chosen abelian
cocycles, factor through \(\Gamma_{S'}=G_{k,S'}\) for any
finite \(S'\supset S\).  Let
\(\mathcal E_{c,S'}=\mathcal Q_c\rtimes\Gamma_{S'}\).
Writing \(\pi_{S'}:G_k\to\Gamma_{S'}\), its actual
pullback isomorphism is
\[
 G_k\times_{\Gamma_{S'}}\mathcal E_{c,S'}
       \xrightarrow{\ \sim\ }\mathcal Q_c\rtimes G_k,
 \qquad (g,(q,\pi_{S'}g))\longmapsto(q,g).
\tag{A.196.17}
\]
Thus a supported lift means a section over \(\Gamma_{S'}\)
whose pullback has the required frame, or equivalently a
geometric cocycle factoring through \(\Gamma_{S'}\).

Fix the literal frame \((a_c,b_c)\).  The set \(Z_c(S')\)
of its cocycle lifts over \(\Gamma_{S'}\) is finite.  To see
this with the action retained, let \(H_c\) be the finite
image of the action on \(\mathcal Q_c\).  Every such cocycle
gives a homomorphism into the finite group
\(\mathcal Q_c\rtimes H_c\).  Its kernel defines a
bounded-degree extension unramified outside \(S'\).
There are finitely many such extensions, and finitely many
maps from each of their finite Galois groups into the
displayed finite group.  This proves the assertion for
literal cocycles.  Abelian discrepancies of representatives
can be removed as in A.196.9, so it covers the class-valued
lifting question too.  Equivalently, the finite-coefficient
first-cohomology finiteness of
\cite[I, Corollary 4.15]{BIR_MilneADT}, applied successively
to the displayed central series, proves the same assertion.

The reduction maps give an actual inverse system of these
finite solution sets.  Nonempty solution sets at unbounded
levels would make every set nonempty.  Every finite initial
compatibility problem can then be solved by reducing a
solution from its largest index.  Compactness of the product
of these finite sets supplies an infinite compatible system.
Its pullback along A.196.17 would contradict A.196.10.
Therefore some level fails, and all higher levels fail by
reduction.  For \(S'=S\), level one has the literal section
\(g\mapsto(x^{a_1(g)},g)\), since \(b_1=0\).
At level two the supported obstruction is precisely
\(\mathfrak o_{S,2}\bmod9\ne0\) from A.195.  This proves
the exact first supported failure claimed in the statement.

\emph{The actual local point and cusp comparison.}
There is one place of \(k\) above three.  Write
\(k_3=\Q_3(\alpha)\), with valuation \(v(\alpha)=1\),
so \(v(3)=3\).  The elementary identities and their depths are
\[
\begin{gathered}
 U-1=\alpha^2(1-\alpha),\qquad
 B-1=-6\alpha^2U,\\
 v(A)=1,\qquad v(B)=0,\qquad
 v(U-1)=2,\qquad v(B-1)=5.
\end{gathered}
\tag{A.196.18}
\]
For \(z\in\alpha\mathcal O_{k_3}\), the three terms of
\((1+z)^3-1=3z+3z^2+z^3\) all have valuation at least
three.  Cubing once more shows
\[
 Y\equiv1\pmod\alpha\quad\Longrightarrow\quad
 v(Y^3-1)\ge3,\qquad v(Y^9-1)\ge6.
\tag{A.196.19}
\]
These bounds follow by the valuation of each term; they do
not require an equality in either minimum.

The actual ordinary-point map at exponent nine is the cover
\[
 A X_0^9+B X_1^9=1,\qquad
              t=A X_0^9,\quad 1-t=B X_1^9.
\tag{A.196.20}
\]
It has no \(k_3\)-point with \(X_0X_1\ne0\).
Indeed \(v(t)\equiv1\pmod9\), so it is nonzero.  If it
were negative, \(v(1-t)=v(t)\) would contradict
\(v(BX_1^9)\equiv0\pmod9\).  If it is positive, then
\(X_1\) is a unit with residue one: the residue field is
\(\F_3\), \(B\equiv1\), and the ninth-power map is the
identity on \(\F_3^\times\).  Equations A.196.18--19 give
\(v(1-BX_1^9)=5\), contradicting \(v(t)\equiv1\pmod9\).

Also \(B\) is not a ninth power in \(k_3\).  This field
has no \(\mu_3\), since its degree over \(\Q_3\) is
three whereas \([\Q_3(\mu_3):\Q_3]=2\).  Thus
\(Y^9=U^3\) would give \(Y^3=U\).  Its residue forces
\(Y\equiv1\pmod\alpha\), and A.196.19 then contradicts
\(v(U-1)=2\).

Finally the functions \(t,1-t\) send the three cuspidal
tangent classes modulo ninth powers to
\((\tau,1)\), \((1,\tau)\), and \((\tau,\tau)\),
respectively.  In the last expression the factor \(-1\)
is itself a ninth power, so it has the displayed trivial
class.  The first line fails because \(B\) is not a ninth
power; the other two fail because the two valuations of
\((A,B)\) are one and zero modulo nine.  These are the
actual abelian Kummer images of the local point and cusp
maps.  They prove the stated scope comparison with the
norm-derived local sections of A.181.

The finite class-three group laws, arithmetic actions,
coefficient maps, root-torsor quotient identities, and all
residue exponents in the proof are checked in
\path{code/two_inert_class_three_lifting_certificate.py},
with retained output and data.  The nonzero global class is
the already proved A.195 class; no finite absolute obstruction
is substituted for it.
\end{proof}
